\documentclass[11pt,a4paper]{article}
\usepackage[margin=23mm]{geometry}
\usepackage{amsmath,amssymb,array,booktabs}
\input{glyphtounicode}
\pdfgentounicode=1
\usepackage[hidelinks]{hyperref}
\hypersetup{hypertexnames=false}
\setlength{\parindent}{0pt}
\setlength{\parskip}{5pt}
\newcommand{\pos}[1]{[#1]_+}
\renewcommand{\labelitemi}{\(\bullet\)}
\begin{document}
\begin{center}
{\Large Selected O-R: causal renewal and a complete analytical cohort}\\[4pt]
15 October 2026 --- mathematical qualification; no traffic experiment
\end{center}

\textbf{Status.} The user selects renewal; RE-O1 closes. Qualify the
previously specified recent-mass O-R formula, retaining R-E and S-E-P.
No stateless replacement, step-dependent reset or new physical formula is
introduced. The bounded recovery candidate/analytical pass is complete at
\textbf{OR-Q1}. Whole-range event qualification, applied MP6 and Step 3
remain open. Parameters and representative speeds remain hypotheses.

\section*{1. The conditional member being qualified}

Keep H1 preparation, cutoff, initially empty T1 and constant advice
$v\in[v_-,v_+]$. For $j=A,E,B$,
\begin{equation}
 P_j(t)=F_j(t-b/v),\quad q_j=P_j-Y_j,\quad N_j=F_j-Y_j,
 \quad \lambda_j=\dot P_j,\quad g_j=\dot Y_j. \tag{1}
\end{equation}
The B lane is independent and unqueued. $B(t)$ denotes instead the
shared A/E budget. With $K=wk_j$,
\begin{equation}
 c(u)=\frac{Ku}{w+u},\quad C(v)=c(v),\quad U(f)=\frac{wf}{K-f},
 \quad 0<R\leq C(v_-),\quad u_R=U(R). \tag{2}
\end{equation}
Use the existing trial work family and sufficient monotonicity domain:
\begin{equation}
 c_A=1,\quad c_E=c_0+\beta\frac{g_A}{C(v)}\frac{\pos{u_A-u_E}}v,
 \quad c_0=1+\delta,\quad
 \eta=c_0-\beta(1+v_+/w)>0. \tag{3}
\end{equation}
Here $\delta,\beta\geq0$, with free marks $u_j=v$ and queued marks
$u_j=U(g_j)$. S-E-P gives a fitting free A offer priority; queued A
uses the earlier equal-work/unused-work template. Its selected-regime
inverse is unique under (3). These remain trial constitutive assumptions,
not a field-calibrated fairness or crossing-speed law.

For positive physical acceleration $a_r$,
\begin{align}
 V(\theta;u)&=\sqrt{(w+u)^2+2wa_r\theta}-w,\nonumber\\
 D(\theta;u)&=\pos{R-c(V(\theta;u))},\quad \theta\geq0,\nonumber\\
 W(u)&=\frac{k_c(u_R)\pos{u_R-u}^2}{2a_r},\quad
 k_c(z)=\frac K{w+z}. \tag{4}
\end{align}
$D$ decreases with age and mark, $0\leq D\leq R$,
$|D_\theta|\leq a_rk_j$, $|D_u|\leq k_j$, and expires within
$\Theta_{\max}=u_R/a_r+u_R^2/(2a_rw)$. Write $W_0=W(0)$.
The age includes backward propagation and is independent of numerical steps.

\newpage
\section*{2. The selected renewal formula and its physical-time update}

Let $Y=Y_E$, $g=g_E$, and $m=m_{\rm ref}=1$ vehicle-equivalent. For
the mass-rank inverse $T(\xi)$ of $Y$, define $a(t)=\pos{Y(t)-m}$ and
\begin{equation}
 p(t)=\frac1m\int_{a(t)}^{Y(t)}
 D(t-T(\xi);u_E(T(\xi)))d\xi,\qquad B(t)=\pos{R-p(t)}. \tag{5}
\end{equation}
This is the previously proposed O-R law. The reference mass is physical,
not a numerical bin or integer admission. Only E creates profiles; A
still affects E admission and current sharing. Traffic mass is never
removed when a profile retires.

The equivalent time-history form is
\begin{equation}
 p(t)=\frac1m\int_0^t D(t-s;u_E(s))g(s)
                  {\bf1}_{\{Y(s)>a(t)\}}ds. \tag{6}
\end{equation}
Flat portions of $Y$ have zero admission measure, so inverse choices and
rank endpoints do not change the integral. The history can be restricted
to both age $\Theta_{\max}$ and the latest admitted mass $m$.

Where derivatives exist, let $T_o=T(Y(t)-m)$ when $Y(t)>m$. Then
\begin{align}
 \dot p={}&\frac{g(t)}mD(0;u_E(t))\nonumber\\
 &+\frac1m\int_{a(t)}^{Y(t)}
 D_\theta(t-T(\xi);u_E(T(\xi)))d\xi\nonumber\\
 &-\frac{g(t)}m{\bf1}_{\{Y(t)>m\}}
                    D(t-T_o;u_E(T_o)). \tag{7}
\end{align}
The three terms are new entry, aging of retained profiles, and retirement
of oldest weight. No division by the oldest discharge rate is needed.
At $Y=m$, use the a.e. version; the cumulative account is continuous.
With $g=0$, only aging remains, so recovery continues without new discharge.

The reading alone generally does not determine future aging or retirement:
recent age/mark information is part of the state. New fast E has zero new
loss but can retire an old slow profile. This is the accepted renewal
approximation to qualify physically, not an exact moving-leader solver.
It does not reset at queue clearance or each numerical step.

\newpage
\section*{3. Causal construction including the moving mass cutoff}

Compare two admitted histories $(g_i,u_i,Y_i)$ with identical past through
$t_0$, $0\leq g_i\leq G:=R/c_0$, and $0\leq u_i\leq v_+$.
For fixed $t>t_0$, put
\begin{equation}
 \epsilon=\sup_{s\leq t}|Y_1(s)-Y_2(s)|
       \leq\int_{t_0}^t|g_1-g_2|ds,\qquad
 |a_1(t)-a_2(t)|\leq\epsilon. \tag{8}
\end{equation}
Use $f_i(s)=D(t-s;u_i(s))$ and the indicator $I_i$ from (6).
The difference with $I_1$ held fixed is bounded by
$R\int|\Delta g|+k_jG\int|\Delta u|$. Where $I_1\ne I_2$,
$|Y_2(s)-a_2(t)|\leq2\epsilon$. The $g_2ds$ mass in that rank band
is at most $4\epsilon$, including plateaus. Thus
\begin{equation}
 |p_1(t)-p_2(t)|\leq\frac1m\left[
 5R\int_{t_0}^t|\Delta g|ds+
 k_jG\int_{t_0}^t|\Delta u|ds\right]. \tag{9}
\end{equation}
This estimate needs no positive lower admission rate or smooth inverse.
It also allows bounded discontinuities in already recorded speed marks.

On a valid selected node-regime chart, (3) gives local uniform Lipschitz
constants $L_g,L_u$ for $g_E$ and its mark as functions of $B$ with
the same arrival data. Since the budget map is 1-Lipschitz,
\begin{equation}
 \|\Delta\mathcal B\|_{\infty,[t_0,t_0+\tau]}
 \leq\frac{\tau}{m}(5RL_g+k_jGL_u)
             \|\Delta B\|_{\infty,[t_0,t_0+\tau]}. \tag{10}
\end{equation}
For sufficiently short $\tau$, this is a contraction on a small valid
chart neighborhood. Bounded account increments preserve that neighborhood
until its next guard. Successive local intervals give causal continuation
while the chart stays valid, and the same estimates give finite-interval
continuous dependence there. The interval length is a proof device.

Equation (10) is not a proof of every endogenous queue switch or a claim
that a free chart is valid when its offer fails. No fallback, numerical
hysteresis or unapproved mode extension is used. The strict transition
checks already established for S-E-P remain separate inputs. The finite
event family below closes its own particular switching sequence.

\newpage
\section*{4. Derived uniform positive service and finite clearance}

The retained E mass is at most $m$ and its rate at most $G=R/c_0$.
Because $D(\theta;u)\leq D(\theta;0)$ decreases with age, its largest
possible integral fills the most recent ages at rate $G$. Therefore
\begin{equation}
 p(t)\leq p_{\max}:=\frac Gm\int_0^{m/G}D(\theta;0)d\theta<R,
 \qquad B(t)\geq B_{\min}:=R-p_{\max}>0. \tag{11}
\end{equation}
The strict inequality follows from $a_r>0$: $D(\theta;0)<R$ at every
positive age. The rearrangement bound is valid even when less than $m$
has passed. It is uniform in time and advice for the fixed parameters.
It is derived, not an added service floor; it can approach zero as recovery
becomes arbitrarily slow. The auxiliary $a_r=0$ limit below is excluded
from this bound and the admitted model.

Let $\overline M_E=P_E(L/v)$, $Q=q_A+q_E$ and $c_{\max}=c_0+\beta$.
The O-R integrand is a subset of the additive one, so Tonelli gives
\begin{equation}
 \int_0^\infty p(t)dt\leq\overline M_E W_0/m. \tag{12}
\end{equation}
For any admissible joined solution, work conservation after $t_*=L/v$
and (11)--(12) give either sufficient bound
\begin{align}
 T_{\rm clear}&\leq t_*+c_{\max}Q(t_*)/B_{\min},\nonumber\\
 T_{\rm clear}&\leq t_*+
       \{c_{\max}Q(t_*)+\overline M_EW_0/m\}/R. \tag{13}
\end{align}
These use available finite E mass before clearance, not assumed final
completion. They do not prove existence of every joined solution.

\textbf{Limits and simpler rivals.} With no E, $p=0$. When free offers
fit from the H1 front, decreasing E offer keeps $p=q_A=q_E=0$ under the
causal free-mode guard. For a constant free A offer $a<R$, (7) gives
$\dot p\leq Rg_E/m$ and $g_E\leq(B-a)/c_0$. The gap $x=B-a$ therefore
satisfies $\dot x\geq-Rx/(mc_0)$ and remains positive at finite times.
Renewal preserves the constant-A protection limit; increasing A arrivals
need a separate argument. For constant queued E rate and mark over the
renewal window,
\begin{equation}
 p=\frac{g_E}{m}W_{m/g_E}(u_E),\quad
 W_h(u):=\int_0^hD(\theta;u)d\theta,\quad
 g_A+c_Eg_E=R-\frac{g_E}{m}W_{m/g_E}(u_E). \tag{14}
\end{equation}
This is the restricted stationary target for a simpler state-dependent
instantaneous rival, not its complete transient selection theorem. O-S,
fixed M1, aggregate/history/scalar and search rivals remain. History must
earn a consequential transient role; no interior optimum is required.

\newpage
\section*{5. A nonempty H1 build-up with A delay: auxiliary prefix}

This is an existence construction, not a chosen traffic scenario, calibration
or optimization result. Fix one advice $v_0$ and the translated eligible
pulse, initially with empty queues/account. Choose a donor-limited H1 front,
so during the pulse
\begin{equation}
 e(t)=e_f e^{-kt},\quad a(t)=s-e(t),\quad k=v_0/h>0,\quad
 s=v_0r<R,\quad \varepsilon=R-s>0. \tag{15}
\end{equation}
Here $a,e$ are \emph{arrival rates}; $h$ is the retained H1 opportunity
length. Choose $c_0>1$, $e_f>0$, $a_f=s-e_f>0$ such that
\begin{equation}
 \varepsilon<\delta e_f/2,\quad
 G_0=(\varepsilon+e_f)/c_0<R,\quad
 d_0=R-G_0>\varepsilon+e_f,\quad \delta=c_0-1. \tag{16}
\end{equation}
Initially $a_f+c_0e_f=R-\varepsilon+\delta e_f>R$, so the strict
S-E-P front activates E while A stays free.

First set $a_r=\beta=0$ solely as a finite-prefix auxiliary limit. Before
the renewal window fills and while A is free and E queued, (5) becomes
\begin{equation}
 g_E=\frac{B-a(t)}{c_0},\quad
 \dot B=-\frac{g_E(R-g_E)}m,\quad
 \dot q_E=e-g_E. \tag{17}
\end{equation}
The queued trace gives $D(0;U(g_E))=R-g_E$. While A is free,
$g_E\leq G_0$ and $B\geq a\geq a_f$. Hence
\begin{equation}
 R-B\geq d_0Y_E/m,\quad
 Y_E\leq m\frac{\varepsilon+e_f}{d_0}<m. \tag{18}
\end{equation}
This rules out a first window fill before A onset, validating the prefix.

Until $B$ reaches $s$, if $e\geq e_f/2$, then
\begin{equation}
 -\dot B\geq\mu:=\frac{e_fd_0}{2mc_0}>0,\quad
 \dot q_E\geq(\delta e_f/2-\varepsilon)/c_0>0. \tag{19}
\end{equation}
Choose $k$ sufficiently small that $k\varepsilon/\mu<\log2$.
Then $B$ reaches $s$ by $\varepsilon/\mu$ while that bound holds and
passes strictly below $s$. Thereafter, while A remains free,
$g_E<e/c_0<e$, so E cannot clear. Since $B$ is nonincreasing and
$a(t)\uparrow s$, a finite first A onset $t_A$ must occur. At that hit,
\begin{equation}
 x=B-a=0,\quad \dot x=-ke(t_A)<0,\qquad
 \dot q_A(t_A+) =a(t_A)-B(t_A)/2=B(t_A)/2>0. \tag{20}
\end{equation}
The first equality uses the pre-onset side; the last uses queued cooperation.

\newpage
\section*{6. Ensuring an E backlog and a nontrivial renewal event}

The construction must exercise renewal, not only its additive initial
prefix. At $k=0$ in the same auxiliary input, A is constant $a_f$ and
the usual barrier keeps $B>a_f$ at every finite time. E stays queued with
\begin{equation}
 \dot q_E\geq\sigma:=(\delta e_f-\varepsilon)/c_0>0,
 \qquad Y_E<m. \tag{21}
\end{equation}
Choose finite $T_b$ with $\sigma T_b>3m$. On this fixed prefix the free-A
gap is strictly positive. Continuous dependence of (17) on the input makes
it possible to choose a sufficiently small \emph{positive} $k$ so A stays
free through $T_b$ and $q_E(T_b)>2m$. Condition (19) can be retained.
The finite A onset from (20) is later, with still larger E backlog.
No constant-input or $k=0$ primary task is adopted.

End the actual finite pulse at $T>t_A$ sufficiently close to $t_A$.
Then $q_A(T)>0$, $q_E(T)>2m$, $Y_E(T)<m$ and $B(T)>0$.
Choose $T$ closer if necessary so
\begin{equation}
 \frac{q_A(T)}{c_0}<m-Y_E(T),\quad
 q_A(T)<c_0q_E(T),\quad
 \frac{R q_A(T)}{mc_0}<B(T)/2. \tag{22}
\end{equation}
These inequalities are possible because $q_A(T)\downarrow0$ as
$T\downarrow t_A$ while the other margins remain positive.

In the auxiliary zero-recovery continuation, while both queues drain and
the window is unfilled, $g_A=B/2$, $g_E=B/(2c_0)$. Thus A clears first,
using at most $q_A(T)/c_0$ E mass, before window fill. Since
$dB/dY_E=-(R-g_E)/m\geq-R/m$, (22) keeps $B\geq B(T)/2$ until
that strict A clearance. Afterwards A has no arrivals and $g_E=B/c_0$.

The first $Y_E=m$ is reached at finite time in this auxiliary E-only
prefix. Before it, $B\geq R(1-Y_E/m)>0$. If $Y_E$ approached $m$
only asymptotically, the additional identity
\begin{equation}
 B=R(1-Y_E/m)+\frac1m\int_0^{Y_E}g_E(T(\xi))d\xi \tag{23}
\end{equation}
would give a strictly positive limiting budget, since some earlier E has
passed at positive rate. That contradicts asymptotic approach. Thus fill
occurs with $g_E>0$ and a positive E backlog. The oldest profile then has
positive loss $R-g_E(0+)>0$. Immediately after fill, the retirement term
in (7) is genuinely nonzero. This finite auxiliary prefix is all that is
needed; its nonrecovering infinite tail is not an admitted model.

\newpage
\section*{7. Positive recovery, positive mismatch and complete clearance}

On every fixed finite interval, (4) approaches $\pos{R-c(u)}$ uniformly
as $a_r\downarrow0$, with error at most $a_rk_j$ times the interval
length. The node inverse is uniformly regular for sufficiently small
$\beta\geq0$ under (3). Equations (9)--(10) give continuous dependence
on these kernels/parameters within valid regimes.

All prefix events used above have strict margins or transverse hits:
E front shortfall, A birth in (20), fixed pulse end with positive queues,
A clearance at positive service, and renewal fill at positive $g_E$.
Between events local continuous dependence applies. At each hit a nonzero
guard derivative bounds the hitting-time change; the compatible new regime
then continues it. Induction through this finite event list preserves
the strict inequalities for sufficiently small \emph{positive} $a_r$ and
\emph{positive} $\beta$. Oldest-profile loss at first retirement also stays
positive. This removes both auxiliary zero parameters, and preserves an
interval with $q_A>0$ and nonzero continuing waiting area.

O-S started from the same inputs agrees through first fill $t_R$. At that
shared state, physical aging and injection agree, but only O-R retires old
weight. Since both budgets are positive,
\begin{equation}
 (\dot B_{O\text{-}R}-\dot B_{O\text{-}S})(t_R+)
   =\frac{g_E(t_R)}mD(t_R;u_E(0+))>0. \tag{23a}
\end{equation}
This is a consequential local difference, not a global trajectory or policy
ordering. It does not establish that either closure is empirically correct.

After the finite pulse there are no new arrivals. Nonnegative queues are
nonincreasing; at most two clearances remain. Bounds (11)--(13) make them
finite; the clearance directions are compatible. The chosen small A tail
makes A clear first with E still positive. E alone then clears at positive
service. Once both queues are zero, no birth is allowed without arrivals;
the fixed retained profiles simply age to zero. This proves a complete
unique regular trajectory for the constructed family, not every possible
switch sequence in the full parameter/action domain.

\textbf{Realizing the inputs with unchanged H1.} Choose $r=s/v_0$,
$e_0+a_0=r$, $a_0=b_0>0$. A receiving-limited initial branch followed
by a strictly donor-limited front is possible when
\begin{equation}
 k_c(v_0)-r<(\gamma-1)e_0,
 \qquad e_f/v_0<\{k_c(v_0)-r\}/(\gamma-1). \tag{24}
\end{equation}
Choose the front progress $z_f$ after that switch. Increasing $h$ while
holding $z_f/h$ fixed retains the same front density and makes $k=v_0/h$
small. Put $\ell=v_0T$, $L_p=\ell+z_f$, and choose finite terminal
length $b>M/k_j$, $M=\ell(r+b_0)$. This changes the common eligibility
delay, not the translated node argument. Initial densities, $v_+r<R$,
$R\leq C(v_-)$ and an adequate independent B lane can all be retained.

The finite strict event sequence persists in a sufficiently small
one-sided neighborhood of $v_0=v_+$, inside the declared action interval.
This is local forward coverage, not permission to shrink that interval
or assume the remaining advices are qualified. Long geometry and very
slow recovery may be needed by this proof. No applied parameter range,
preferred advice or interior optimum is established.

\newpage
\section*{8. Exact feasibility, accounting and the review checkpoint}

An exact rational point verifies that the static inequalities are nonempty:
\begin{equation}
 R=m=1,\quad c_0=11/10,\quad e_f=1/8,\quad
 \varepsilon=1/512,\quad s=511/512. \tag{25}
\end{equation}
Then
\begin{align}
 d_0&=2491/2816,\quad \mu=12455/247808,\quad
 \varepsilon/\mu=484/12455,\nonumber\\
 \sigma&=27/2816,\qquad
 (\varepsilon+e_f)/d_0=715/4982<1. \tag{26}
\end{align}
Take $v_-=1,v_0=v_+=2,w=2,k_j=3$, $a_0=b_0=1/8$,
$e_0=383/1024$, $r=511/1024$, $\gamma=5$, $R_B=1$.
The fast initial receiving branch is strict; the slower one is donor-limited.
A donor-limited fast front with density $1/16$ is obtained from
\begin{equation}
 \frac{z_f}{h}=5\log(5632/5125)+\log(1025/256). \tag{27}
\end{equation}
$k$ is subsequently chosen sufficiently small by the finite-prefix argument;
no numerical traffic trajectory, terminal size or positive recovery/mismatch
parameter set is selected. These values are algebraic feasibility evidence.

For the resulting admissible family, T1's storage proof is unchanged:
$M<k_jb$ and $\dot F_j\leq C(v)$ keep entrance vacancy positive.
Early-merged mass enters A once. With remaining upstream occupied length
$s_{\rm occ}$ and initial total density $\rho_0$,
\begin{equation}
 s_{\rm occ}\rho_0+\sum_jN_j+\sum_jY_j=M,\qquad
 J(v)=\frac{M(L-\ell/2)}v+\int_0^\infty(q_A+q_E)dt. \tag{28}
\end{equation}
Retiring recovery information is not removing traffic. Post-line influence
is not post-line vehicle-time. Known empty memory and initial continuing
symmetry retain conditional total-plus-ending-statistic sufficiency for
uniquely qualified joined cases. Neither full-lane-state necessity nor a
positive information/action/gradient result follows.

\textbf{Source boundary.} The isolated point motivation is reused from
Leclercq et al. (2016), Section 2; O-R and all proofs here are our adaptation
and deductions. Jin's LTM boundary/node context does not make this a general
invariant junction. The representative speed, renewal mass, maneuver/recovery
partition and common-advice response still require physical qualification.

\textbf{Checkpoint OR-Q1.} The bounded J3 candidate-and-analytical-check
pass is complete. O-R, R-E and S-E-P remain selected. The regular cohort
advances mathematical coexistence, not applied MP6. Stop for review before
J4's original-action-range, provenance and information/sensitivity/native
qualification. Step 3 remains open, Step 4 unopened. No simulator,
installation, traffic experiment, numerical gradient or optimization ran.
\end{document}
